\subsection{ANTI-INVARIANT ELEMENTS}

We retain the assumptions and notation of Th. 3, and assume that K is of characteristic 0. An element z of S is said to be anti-invariant under G if

\[
g (z) = (\det g) ^ {- 1} z
\]

for all \(g\in G\)

Let H be the set of pseudo-reflections belonging to G and distinct from 1. For all \(g \in H\) , there exist \(e_{g} \in V\) and \(f_{g} \in V^{*}\) such that

\[
g (x) = x + f _ {g} (x) e _ {g} \quad \text {   for   all   } x \in V.
\]

PROPOSITION 5. (i) Let \(\mathrm{D}\) be the element \(\prod_{g \in \mathrm{H}} e_g\) of \(\mathrm{S}\) . The elements of \(\mathrm{S}\) which are anti-invariant under \(\mathrm{G}\) are the elements of \(\mathrm{RD}\) .

\begin{enumerate}
\def\labelenumi{(\roman{enumi})}
\setcounter{enumi}{1}
\tightlist
\item
  Identify S with the polynomial algebra \(\mathrm{K}[\mathrm{X}_1,\dots ,\mathrm{X}_l]\) by choosing a basis \((\mathrm{X}_1,\dots ,\mathrm{X}_l)\) of V, and let \((\mathrm{P}_1,\dots ,\mathrm{P}_l)\) be algebraically independent homogeneous elements of S generating the algebra R (Th. 3). Then the jacobian \(\mathrm{J} = \det \left(\frac{\partial\mathrm{P}_t}{\partial\mathrm{X}_j}\right)\) is of the form \(\lambda D\) , where \(\lambda \in K^*\) .
\end{enumerate}

\begin{enumerate}
\def\labelenumi{\alph{enumi})}
\tightlist
\item
  With the notation in (ii), we have
\end{enumerate}

\[
d \mathrm{P} _ {1} \wedge d \mathrm{P} _ {2} \wedge \dots \wedge d \mathrm{P} _ {l} = \mathrm{J} d \mathrm{X} _ {1} \wedge d \mathrm{X} _ {2} \wedge \dots \wedge d \mathrm{X} _ {l},
\]

so, for all \(g \in G\) ,

\[
\begin{array}{c} g (\mathrm{J}) (\det g)   d \mathrm{X} _ {1} \wedge \ldots \wedge d \mathrm{X} _ {l} = g (\mathrm{J})   d (g \mathrm{X} _ {1}) \wedge \ldots \wedge d (g \mathrm{X} _ {l}) \\ = g (d \mathrm{P} _ {1} \wedge \ldots \wedge d \mathrm{P} _ {l}) = d \mathrm{P} _ {1} \wedge \ldots \wedge d \mathrm{P} _ {l} = \mathrm{J}   d \mathrm{X} _ {1} \wedge \ldots \wedge d \mathrm{X} _ {l}, \end{array}
\]

hence J is anti-invariant under G. On the other hand, the field of fractions N of S is a Galois extension of the field of fractions E of R (no. 2); if \(\Delta\) is a derivation of E with values in an extension field \(\Omega\) of N, \(\Delta\) extends to a derivation of N with values in \(\Omega\) (Algebra, Chap. V, § 16, no. 4, Th. 3); since the \(P_{i}\) are algebraically independent, it follows that \(dP_{1} \wedge \ldots \wedge dP_{l} \neq 0\) , hence \(J \neq 0\) .

\begin{enumerate}
\def\labelenumi{\alph{enumi})}
\setcounter{enumi}{1}
\item
  Let \(z\) be an element of \(S\) anti-invariant under \(G\) . We show that \(z\) is divisible by \(D\) in \(S\) . Let \(a\) be a non-zero vector in \(V\) . The elements of \(G\) that leave \(Ka\) stable leave stable a complementary hyperplane \(L\) (Appendix, Prop. 2); an element of \(G\) leaving \(Ka\) stable is 1 or a pseudo-reflection with vector \(a\) if and only if it induces 1 on \(L\) ; the pseudo-reflections with vector \(a\) belonging to \(G\) thus constitute, together with 1, a cyclic subgroup \(G'\) of \(G\) ; let \(t\) be its order. There exists a basis \((X_1, \ldots, X_l)\) of \(V\) such that \(a = X_1\) , \(X_2 \in L, \ldots, X_l \in L\) , and \(z\) can be identified with a polynomial \(P(X_1, \ldots, X_l)\) with coefficients in \(K\) . From \(g(z) = (\det g)^{-1}z\) for \(g \in G'\) , we see that \(X_1\) only appears in \(P(X_1, \ldots, X_l)\) with exponents congruent to -1 modulo \(t\) . Thus, \(P(X_1, \ldots, X_l)\) is divisible by \(X_1^{t-1} = a^{t-1}\) . Now \(D\) is, up to a scalar factor, the product of the \(a^{t-1}\) for those \(a \in V\) such that \(t > 1\) , and these elements of \(S\) are mutually coprime. Since \(S\) is factorial, \(z\) is divisible by \(D\) .
\item
  By a) and b), J is divisible by D in S. Now
\end{enumerate}

\[
\deg \mathrm{J} = \sum_ {i = 1} ^ {l} (k _ {i} - 1) = \operatorname{Card} (\mathrm{H})
\]

(Prop. 3), so \(\deg J = \deg D\) , hence \(J = \lambda D\) with \(\lambda \in K\) . Since \(J \neq 0\) , \(\lambda \in K^*\) . This proves (ii).

\begin{enumerate}
\def\labelenumi{\alph{enumi})}
\setcounter{enumi}{3}
\tightlist
\item
  Parts a) and c) of the proof show that D is anti-invariant under G. Next, if \(y \in R\) , it is clear that yD is anti-invariant under G. Finally, if \(z \in S\) is anti-invariant under G, we have seen in b) that there exists \(y \in S\) such that z = yD. Since S is integral, \(y \in R\) . This proves (i).
\end{enumerate}
% semantic-fragments: legacy-input-seam
\relax{}
